In a toy ground-aerial open-vocabulary search task where only the message channel changes, templated natural-language messages did not beat fixed-schema symbolic messages on time-to-find; pooled over query types they were slightly slower, with a small gain in success. The gain appears where the hypothesis placed it, on attribute and relational queries, and disappears or reverses for exact and synonym queries. The decomposition locates the causes: attribute and relation content helps in either format, and a room-name reference costs a couple of steps compared with coordinates. Under message corruption both channels degrade and symbolic messaging is not the more robust one under our corruption models. Any communication is worth far more than the choice of format.

For system builders the practical reading is that a language channel earns its place through what it can say that the schema designer did not anticipate, not through being language; where a field can be added, a field does the same job at half the tokens, and precise spatial references should be kept alongside the sentence. Whether real language models change this picture, through inference or through error, is outside what this study can say and is the natural next experiment.
